\subsection{Encoder and Competency-Type Comparisons}
\label{sec:external-benchmark-results}

XLM-R-large and ESCOXLM-R achieved mean typed exact-span micro-F1 scores of 0.522 and 0.538, respectively (Table~\ref{tab:external-chinese}, panel A). Scores were recomputed from all six archived prediction files against the same human reference set. The mean difference was 0.016, with a paired 95\% interval of $[-0.005, 0.035]$ (Table~\ref{tab:core-paired}). The interval includes zero and does not establish an advantage for ESCO initialization under this protocol.

Four-type macro-F1 was 0.607 for XLM-R-large and 0.622 for ESCOXLM-R. Averaging over K/S/T gave 0.542 and 0.562, respectively. These sensitivity summaries exclude the sparsely represented L category from the macro average without changing the four-type prediction task.

Qwen adaptation improved all three K/S/T categories (Table~\ref{tab:external-chinese}, panel B; Figure~\ref{fig:qwen-diagnostics}). The largest observed increase among these categories was for T, from 0.380 to 0.638. Mean S and K scores increased by 0.148 and 0.210, respectively. Among the three Silver-trained configurations in Table~\ref{tab:external-chinese}, panel B, Qwen had the highest observed mean scores for S and T, while ESCOXLM-R had the highest K score. These rankings are descriptive because prediction heads, training budgets, and inference procedures differ across architectures.

\begin{table}[!htbp]
\centering\small
\caption{Supervised adaptation of multilingual encoders on Gold150 (0--1), mean $\pm$ sample SD across three seeds.}
\label{tab:external-chinese}
\begin{tabularx}{\linewidth}{@{}Lrrrr@{}}
\toprule
Encoder & Typed exact & Typed relaxed & Boundary & Macro-F1 \\
\midrule
XLM-R-large & $0.522\pm0.022$ & $0.650\pm0.021$ & $0.571\pm0.022$ & $0.607\pm0.016$ \\
ESCOXLM-R & $0.538\pm0.021$ & $0.662\pm0.022$ & $0.582\pm0.024$ & $0.622\pm0.010$ \\
\bottomrule
\end{tabularx}
\par\smallskip\RaggedRight\footnotesize B2 training/development: 2,150/169 records. Fresh linear LSKT head; official \texttt{cnss-lskt-1.2.0} scorer; 150 human-reference records. Boundary is exact span matching with types collapsed; macro-F1 averages L/K/S/T. Gold150 is not a new blind test and contains only two L spans. All three seeds (42/43/44) are included; per-seed scores and selected epochs are retained in the model audit archive.
\end{table}

\begin{table}[!htbp]
\centering\small
\caption{Common-protocol typed exact token-span micro-F1 (0--1), mean $\pm$ sample SD across three seeds.}
\label{tab:external-common}
\begin{tabularx}{\linewidth}{@{}L P{.14\linewidth}P{.25\linewidth}rr@{}}
\toprule
Dataset & Stream & Train/dev/test & XLM-R-large & ESCOXLM-R \\
\midrule
SkillSpan & Skill & 4,800/3,174/3,569 & $0.567\pm0.012$ & $0.577\pm0.006$ \\
SkillSpan & Knowledge & 4,800/3,174/3,569 & $0.698\pm0.021$ & $0.672\pm0.052$ \\
Kompetencer & Skill & 778/346/262 & $0.545\pm0.028$ & $0.527\pm0.004$ \\
Kompetencer & Knowledge & 778/346/262 & $0.618\pm0.035$ & $0.593\pm0.044$ \\
Gnehm ICT & ICT & 19,889/2,332/2,557 & $0.891\pm0.005$ & $0.867\pm0.032$ \\
Green & Native types & 8,669/964/335 & $0.493\pm0.009$ & $0.492\pm0.007$ \\
Sayfullina & Skill & 3,705/1,855/1,851 & $0.920\pm0.004$ & $0.925\pm0.005$ \\
FIJO & Native types & 399/50/50 & $0.343\pm0.019$ & $0.339\pm0.007$ \\
\bottomrule
\end{tabularx}
\par\smallskip\RaggedRight\footnotesize Both encoders use fresh linear BIO heads, 20 epochs, and development-set checkpoint selection. SkillSpan and Kompetencer have separate skill and knowledge streams; Green and FIJO retain their entity types. German data use the public ICT release. The shared protocol differs from the original CRF and NNOSE methods.
\end{table}

